\clearpage
\VerdierPhysicalPageAnchor{0195}%
\begingroup\small\setlength{\parskip}{0pt}\setlength{\emergencystretch}{1.5em}
\setlength{\abovedisplayskip}{4pt plus 2pt minus 2pt}%
\setlength{\belowdisplayskip}{4pt plus 2pt minus 2pt}%

\VerdierSemanticAnchor{ch3d-2-3-1-proof-cont}
ここで $X_n$ は $\operatorname{comp}^{-}(\mathcal A)$ の対象、$u_n$ は高さ
$-n$ の初等射であり、次の図式はホモトピーを除いて可換である：

\VerdierSemanticAnchor{ch3d-2-3-1-limits-square}
\begin{center}
\begin{tikzcd}[column sep=5.5em,row sep=3.0em]
\displaystyle\varprojlim_n X_n \arrow[r] \arrow[d]
  & X \arrow[d,"m"] \\
\displaystyle\varinjlim_n X_n \arrow[r,"\sim"'] & Y
\end{tikzcd}
\end{center}

\VerdierSemanticAnchor{ch3d-2-3-1-normalization}
したがって、命題を証明するには $X=\varprojlim_n X_n$ と仮定してよく、
さらに $n$ が十分小さいとき、$X_n=Y$ かつ $u_n=\operatorname{id}$ である。
$n$ に関する帰納法により、次を構成する：

\VerdierSemanticAnchor{ch3d-2-3-1-induction-data}
\noindent a) 初等射の列
\[
v_n\colon L_n^\bullet\longrightarrow L_{n-1}^\bullet,
\qquad n\in\mathbf Z,\qquad
v_n\text{ という初等射の高さは }-n.
\]

\noindent b) 射の列
\[
w_n\colon L_n^\bullet\longrightarrow X_n
\]
で、次を満たすもの：

\VerdierSemanticAnchor{ch3d-2-3-1-induction-condition-1}
\noindent 1) 各 $n\in\mathbf Z$ に対して、$L_n^\bullet$ の次数 $i$ の成分は、
$i\geq -n$ なら $M$ に属し、次数 $i$ の成分は、$i<-n$ なら $M^\wedge$ に属する。

\VerdierSemanticAnchor{ch3d-2-3-1-induction-square}
\noindent 2) 次の図式が可換である：
\begin{center}
\begin{tikzcd}[column sep=5.0em,row sep=2.8em]
L_n^\bullet \arrow[r,"w_n"] \arrow[d,"v_n"']
  & X_n \arrow[d,"u_n"] \\
L_{n-1}^\bullet \arrow[r,"w_{n-1}"'] & X_{n-1}
\end{tikzcd}
\end{center}
これは各 $n\in\mathbf Z$ に対して成り立つ。

\VerdierSemanticAnchor{ch3d-2-3-1-induction-condition-3}
\noindent 3) 対象 $L_n^\bullet$ は $Y$ に等しく、射 $v_n$ と $w_n$ は恒等射である。
これは $n$ が十分小さいときに成り立つ。
\endgroup
